% ═══════════════════════════════════════════════════════════════════
%  Rapport de Stage – Sécurisation des LLM – CMR 2026
%  Auteur : Omar Chekroun
%  Compilé sur Overleaf (pdfLaTeX)
% ═══════════════════════════════════════════════════════════════════
\documentclass[12pt, a4paper]{report}

% ── Encodage & langue ───────────────────────────────────────────────
\usepackage[utf8]{inputenc}
\usepackage[T1]{fontenc}
\usepackage[french]{babel}
\usepackage{lmodern}

% ── Mise en page ────────────────────────────────────────────────────
\usepackage[
  a4paper,
  left=2.5cm, right=2.5cm,
  top=2.5cm,  bottom=2.5cm,
  headheight=1.8cm
]{geometry}

% ── Graphiques ──────────────────────────────────────────────────────
\usepackage{graphicx}
\usepackage[export]{adjustbox}
\usepackage{float}
\usepackage{tabularx}

% ── Couleurs ────────────────────────────────────────────────────────
\usepackage[dvipsnames]{xcolor}
\definecolor{primary}{RGB}{0,51,102}
\definecolor{accent}{RGB}{0,102,179}
\definecolor{lightgray}{RGB}{245,247,250}

% ── En-têtes & pieds de page ────────────────────────────────────────
\usepackage{fancyhdr}
\pagestyle{fancy}
\fancyhf{}

% Définition du contenu de l'en-tête (utilisé par fancy et plain)
\newcommand{\myheader}{%
  \includegraphics[height=2 cm, valign=c]{logo_cmr.jpg}
  \hspace{10pt}%
  \color{primary}\small\textbf{\nouppercase{\leftmark}}%
}

% Application au style FANCY (pages normales)
\fancyhead[L]{\myheader}
\fancyfoot[C]{\color{primary}\small\thepage}
\renewcommand{\headrulewidth}{0.5pt}
\renewcommand{\headrule}{{\color{primary}\hrule width\headwidth height 0.5pt}}

% Redéfinition du style PLAIN (pages de début de chapitre)
\fancypagestyle{plain}{%
  \fancyhf{}
  \fancyhead[L]{\myheader} % On remet l'en-tête ici aussi !
  \fancyfoot[C]{\color{primary}\small\thepage}
  \renewcommand{\headrulewidth}{0.5pt}
  \renewcommand{\headrule}{{\color{primary}\hrule width\headwidth height 0.5pt}}
}

% ── Titres ──────────────────────────────────────────────────────────
\usepackage{titlesec}

\titleformat{\chapter}[display]
  {\color{primary}\normalfont\huge\bfseries}
  {\color{accent}\chaptertitlename\ \thechapter}
  {10pt}
  {\color{primary}\Huge}
  [{\color{accent}\titlerule[0.6pt]}]

\titleformat{\section}
  {\color{primary}\normalfont\Large\bfseries}
  {\thesection}{1em}{}
  [{\color{accent}\titlerule[0.4pt]}]

\titleformat{\subsection}
  {\color{primary}\normalfont\large\bfseries}
  {\thesubsection}{1em}{}

\titleformat{\subsubsection}
  {\color{accent}\normalfont\normalsize\bfseries}
  {\thesubsubsection}{1em}{}

\titlespacing*{\chapter}{0pt}{0pt}{1.5em}
\titlespacing*{\section}{0pt}{1.5em}{0.8em}
\titlespacing*{\subsection}{0pt}{1.2em}{0.5em}
% FIX #8 – espacement subsubsection ajouté
\titlespacing*{\subsubsection}{0pt}{1.0em}{0.4em}

% ── Listes ──────────────────────────────────────────────────────────
\usepackage{enumitem}
\setlist[itemize]{
  leftmargin=*, itemsep=0.2cm, topsep=0.2cm,
  label={\color{primary}\textbullet}
}
\setlist[enumerate]{leftmargin=*, itemsep=0.2cm, topsep=0.2cm}

% ── Tableaux ────────────────────────────────────────────────────────
\usepackage{booktabs}
\usepackage{tabularx}
\usepackage{array}
\usepackage{colortbl}

% ── Encadrés colorés ────────────────────────────────────────────────
\usepackage{tcolorbox}
\tcbuselibrary{skins, breakable}

\tcbset{
  cmrbox/.style={
    enhanced, breakable,
    colback=lightgray, colframe=primary,
    fonttitle=\bfseries\color{white}, coltitle=white,
    attach boxed title to top left={yshift=-2mm, xshift=4mm},
    boxed title style={colback=primary, rounded corners},
    arc=4pt, boxrule=0.8pt,
    left=6pt, right=6pt, top=8pt, bottom=6pt
  }
}

% ── Micro-typographie & espacement ──────────────────────────────────
\usepackage{microtype}
\usepackage{setspace}
\setstretch{1.2}
\usepackage{parskip}
\setlength{\parskip}{0.5em}

% FIX #4 – hyperref chargé en dernier pour éviter les conflits
\usepackage{hyperref}
\hypersetup{
  colorlinks      = true,
  linkcolor       = primary,
  urlcolor        = accent,
  citecolor       = accent,
  pdfauthor       = {Omar Chekroun},
  pdftitle        = {Rapport de Stage – Sécurisation des LLM – CMR 2026},
  pdfsubject      = {Stage – Direction de la Transformation Digitale},
  bookmarksnumbered = true,
  bookmarksopen   = true
}

% ════════════════════════════════════════════════════════════════════
%                        DÉBUT DU DOCUMENT
% ════════════════════════════════════════════════════════════════════
\begin{document}

% ────────────────────────────────────────────────────────────────────
%  PAGE DE TITRE
% ────────────────────────────────────────────────────────────────────
\begin{titlepage}
  \thispagestyle{empty}
  \centering

  \includegraphics[width=0.55\textwidth]{cmr.png}\par
  \vspace{0.6cm}

  {\color{primary}\rule{\linewidth}{1.5pt}}\par
  \vspace{0.4cm}

  {\color{primary}\Huge\bfseries Rapport de Stage\par}
  \vspace{0.2cm}
  {\color{accent}\Large Direction de la Transformation Digitale\par}
  \vspace{0.4cm}

  {\color{primary}\rule{\linewidth}{1.5pt}}\par
  \vspace{0.7cm}

  {\Large\bfseries Caisse Marocaine des Retraites (CMR)\par}
  \vspace{0.8cm}

  \begin{tcolorbox}[
    enhanced,
    colback=lightgray, colframe=primary,
    arc=6pt, boxrule=1pt,
    left=10pt, right=10pt, top=8pt, bottom=8pt
  ]
    \centering
    {\color{primary}\large\bfseries Sujet du rapport :\par}
    \vspace{0.3cm}
    {\large\itshape
      Sécurisation des Modèles de Langage de Grande Taille (LLM)\\
      en Environnement d'Administration Publique\par}
    \vspace{0.2cm}
    {\normalsize\color{accent}
      Approches Cloud Public (Pseudonymisation) et On-Premise (Souverain)\par}
  \end{tcolorbox}

  \vspace{0.7cm}

 \begin{tabular}{>{\color{primary}\bfseries}r @{\hspace{0.6em}:\hspace{0.8em}} l}
Durée
  & 5 janvier -- 5 mars 2026 (2 mois) \\

Stagiaire
  & Omar \textsc{Chekroun} \\

Établissement
  & Faculté des Sciences de Rabat (FSR) \\

Master
  & Informatique, Gouvernance et Transformation Digitale \\

Encadrant CMR
  & M.~Anass \textsc{El~Alaoui El~Bahi} \\
  & \small\itshape Digital Transformation Manager \\
\end{tabular}

  \vfill
  {\small\color{primary}
    Caisse Marocaine des Retraites \textbullet{}
    Direction de la Transformation Digitale \textbullet{}
    Rabat, 2026}
\end{titlepage}

% ────────────────────────────────────────────────────────────────────
%  REMERCIEMENTS
% ────────────────────────────────────────────────────────────────────
\newpage
\thispagestyle{plain}
% FIX #3 – en-tête correct pour chapitres non numérotés
\markboth{Remerciements}{Remerciements}
\chapter*{Remerciements}
\addcontentsline{toc}{chapter}{Remerciements}

Je tiens à exprimer ma profonde gratitude à toutes les personnes qui ont
contribué à la réussite de ce stage et à l'élaboration du présent rapport.

En premier lieu à mon encadrant, \textbf{M.~Anass El~Alaoui El~Bahi}, Digital Transformation Manager à la Caisse Marocaine des Retraites. Bien plus qu'un simple superviseur, il a été pour moi un véritable mentor. Je le remercie pour sa disponibilité sans faille, sa bienveillance constante et la confiance qu'il m'a témoignée en me confiant un sujet aussi stratégique. Nos échanges passionnants et ses conseils avisés ont été le moteur de ma réflexion et ont largement contribué à l'épanouissement de mes compétences durant ces deux mois.

Je remercie également l'ensemble de l'équipe de la \textbf{Direction de la Transformation Digitale} de la CMR. Merci pour votre accueil chaleureux, votre soutien quotidien et pour m'avoir permis de m'immerger au cœur des enjeux technologiques de l'institution.

Un remerciement particulier va à mes enseignants de la \textbf{Faculté des Sciences de Rabat}, qui m'ont donné les clés académiques nécessaires pour aborder ce défi avec rigueur.

Enfin, je remercie ma famille et mes proches pour leur soutien indéfectible, qui m'accompagne dans chaque étape de mon parcours.

\vspace{1cm}
\hfill\textit{Omar Chekroun, Rabat, mars 2026}

% ────────────────────────────────────────────────────────────────────
%  RÉSUMÉ & ABSTRACT
% ────────────────────────────────────────────────────────────────────
\newpage
\thispagestyle{plain}
% FIX #3 – en-tête correct
\markboth{Résumé et Abstract}{Résumé et Abstract}
\chapter*{Résumé}
\addcontentsline{toc}{chapter}{Résumé et Abstract}

\textbf{Contexte.} 
L'essor des Modèles de Langage de Grande Taille (LLM) transforme les processus administratifs. Toutefois, pour une institution comme la CMR, l'usage de solutions Cloud publiques soulève des risques majeurs de fuite de données personnelles et de non-conformité réglementaire.

\medskip
\textbf{Problématique.} 
Comment tirer profit de la puissance des LLM tout en respectant strictement la Loi 09-08 (protection des données) et la Loi 05-20 (cybersécurité), sachant que les données des retraités sont hautement sensibles ?

\medskip
\textbf{Approche.} 
Ce rapport détaille la conception d'une architecture de sécurité hybride. D'une part, une \textbf{approche Cloud sécurisée} via un proxy de pseudonymisation utilisant le NER (Named Entity Recognition) et des Regex pour intercepter les PII avant leur sortie du réseau. D'autre part, une \textbf{approche On-Premise} basée sur le déploiement local de modèles open-source (type Mistral ou Llama 3) via des environnements conteneurisés, garantissant une souveraineté totale.

\medskip
\textbf{Résultats.} 
L'étude aboutit à une matrice de décision comparant performance, coût et niveau de risque. Nous préconisons une stratégie "Secure-by-Design" incluant un système de filtrage des prompts, une gestion granulaire des accès (RBAC) et une journalisation exhaustive pour l'auditabilité, alignée sur les exigences de la DGSSI.

\medskip
\textbf{Mots-clés :} LLM, Cybersécurité, Pseudonymisation, NER, Souveraineté, On-Premise, Loi 09-08, Loi 05-20, CMR.

\vspace{1.2cm}

% FIX #6 – thispagestyle ajouté sur la page Abstract
\newpage
\thispagestyle{plain}
% FIX #2 – addcontentsline ajouté pour l'Abstract
\chapter*{Abstract}
\addcontentsline{toc}{chapter}{Abstract}

\textbf{Context.} 
The rise of Large Language Models (LLMs) is redefining administrative workflows. However, for institutions like the CMR, relying on public Cloud solutions introduces significant risks regarding personal data leakage and regulatory non-compliance.

\medskip
\textbf{Problem Statement.} 
How can the CMR leverage LLM capabilities while strictly adhering to Moroccan Laws 09-08 (data protection) and 05-20 (cybersecurity), given the highly sensitive nature of pension-related data?

\medskip
\textbf{Approach.} 
This report outlines a hybrid security architecture. It explores a \textbf{secured Cloud approach} using a pseudonymization proxy powered by Named Entity Recognition (NER) and Regex to scrub PII before data leaves the internal network. Simultaneously, it evaluates an \textbf{On-Premise approach} involving the local deployment of open-source models (e.g., Mistral, Llama 3) to ensure absolute data sovereignty.

\medskip
\textbf{Outcomes.} 
The study provides a decision matrix balancing performance, cost, and risk. We recommend a "Secure-by-Design" strategy featuring prompt filtering, Role-Based Access Control (RBAC), and comprehensive audit logging, fully aligned with DGSSI requirements.

\medskip
\textbf{Keywords:} LLM, Cybersecurity, Pseudonymization, NER, Sovereignty, On-Premise, Law 09-08, Law 05-20, CMR.

% ────────────────────────────────────────────────────────────────────
%  TABLE DES MATIÈRES
% ────────────────────────────────────────────────────────────────────
\newpage
\thispagestyle{plain}
\tableofcontents
\newpage

% ════════════════════════════════════════════════════════════════════
%  CHAPITRE 1 – PRÉSENTATION DE LA CMR
% ════════════════════════════════════════════════════════════════════
\chapter{Présentation de la Caisse Marocaine des Retraites}

% FIX #7 – width réduit à 0.95 pour éviter débordement avec \fbox
\begin{figure}[H]
  \centering
  \fbox{\includegraphics[
    width=0.95\textwidth,
    height=8cm,
    keepaspectratio
  ]{cmr_building.jpeg}}
  \caption{Siège social de la Caisse Marocaine des Retraites, Rabat}
  \label{fig:cmr_building}
\end{figure}

La Caisse Marocaine des Retraites (CMR) est un établissement public à
caractère administratif, doté de la personnalité morale et de l'autonomie
financière. Placée sous la tutelle du Ministère de l'Économie et des
Finances, elle constitue un acteur central du système national de retraite
et de prévoyance sociale au Maroc.

\section{Historique et missions}

Fondée en \textbf{1939}, la CMR a pour mission principale de gérer les
régimes de retraite obligatoires des fonctionnaires et agents publics
marocains. Elle assure la liquidation, le paiement et le contrôle des
prestations, ainsi que le suivi des cotisations pour l'ensemble de ses
adhérents.

Au fil des décennies, la CMR s'est positionnée comme un pilier de la
protection sociale au Maroc, gérant des millions de dossiers de retraités
et de fonctionnaires actifs, ce qui en fait l'une des institutions
détenant les bases de données personnelles les plus volumineuses et
les plus sensibles du pays.

\section{Régimes de retraite gérés}

La CMR administre quatre régimes de retraite distincts :

\begin{itemize}
  \item \textbf{Régime des Pensions Civiles (RPC)} : destiné aux
        fonctionnaires civils de l'État ;
  \item \textbf{Régime des Pensions Militaires (RPM)} : couvrant le
        personnel des Forces Armées Royales ;
  \item \textbf{Régime Collectif d'Allocation de Retraite (RCAR)} :
        pour les agents non titulaires de l'État et des établissements
        publics ;
  \item \textbf{Régime des agents temporaires, occasionnels et
        contractuels} de l'État et des collectivités territoriales.
\end{itemize}

\section{Direction de la Transformation Digitale}

C'est dans le cadre de la \textbf{Direction de la Transformation
Digitale} que s'inscrit le présent stage, sous la supervision de
\textbf{M.~Anass El~Alaoui El~Bahi}, Digital Transformation Manager.
Cette direction est notamment chargée de :

\begin{itemize}
  \item Définir et piloter la stratégie de transformation numérique
        de l'institution ;
  \item Intégrer et sécuriser l'adoption de nouvelles technologies,
        dont l'intelligence artificielle et les solutions cloud ;
  \item Assurer la protection des données sensibles traitées par la CMR,
        conformément aux obligations légales nationales ;
  \item Veiller à la conformité réglementaire en matière de sécurité
        de l'information (Loi~05-20, DGSSI) et de protection des données
        (Loi~09-08, CNDP) ;
  \item Diffuser une culture de sécurité et d'innovation au sein des
        équipes métiers et techniques.
\end{itemize}

La CMR se positionne ainsi non seulement comme un organisme de gestion
administrative, mais également comme un acteur engagé dans la
\textbf{résilience numérique} et la \textbf{transformation digitale
sécurisée}, dans le respect des exigences réglementaires nationales
et des standards internationaux.

% ════════════════════════════════════════════════════════════════════
%  CHAPITRE 2 – OBJECTIFS DU STAGE
% ════════════════════════════════════════════════════════════════════
\chapter{Objectifs du Stage}

\begin{figure}[H]
  \centering
  \includegraphics[width=0.65\textwidth]{obj.png}
  \caption{Axes stratégiques de la sécurisation des LLM à la CMR}
  \label{fig:objectifs}
\end{figure}

% FIX #1 – 2025 corrigé en 2026
Ce stage s'est déroulé du \textbf{5 janvier au 5 mars 2026} au sein de la
Direction de la Transformation Digitale de la CMR, sous la supervision de
\textbf{M.~Anass El~Alaoui El~Bahi}. Il s'inscrit dans une logique de
professionnalisation et de contribution à un défi technologique
d'actualité : la \textbf{sécurisation des LLM en environnement
d'administration publique}.

\section{Analyse des risques LLM en contexte public}

Le premier objectif consistait à analyser les risques inhérents à
l'utilisation de LLM dans un contexte d'administration publique :

\begin{itemize}
  \item Identification des vecteurs d'attaque spécifiques aux LLM
        (\textit{prompt injection}, fuite de données, hallucinations,
        \textit{Shadow AI}) ;
  \item Évaluation de la sensibilité des données traitées par la CMR
        (données personnelles de retraités, données financières,
        informations stratégiques) ;
  \item Comparaison des risques entre approche cloud et on-premise .
\end{itemize}

\section{Étude du cadre juridique marocain}

Un volet essentiel du stage portait sur l'analyse du cadre normatif et
réglementaire applicable à l'usage des LLM dans une institution publique :

\begin{itemize}
  \item \textbf{Loi 05-20} relative à la cybersécurité : obligations
        des entités publiques, rôle de la DGSSI et exigences de
        sécurisation des systèmes d'information critiques ;
  \item \textbf{Loi 09-08} relative à la protection des données à
        caractère personnel : obligations envers la CNDP, consentement,
        minimisation des données et encadrement des transferts hors
        territoire ;
  \item Statut de la CMR en tant qu'\textbf{institution publique gérant
        des données critiques} et contraintes associées ;
  \item Exigences de \textbf{localisation des données} (\textit{data
        residency}) dans un cloud souverain conforme.
\end{itemize}

\section{Exploration technologique : deux approches}

Le stage visait à explorer et comparer deux stratégies de déploiement
des LLM :

\begin{itemize}
  \item \textbf{Approche Cloud (Passerelle)} : utilisation de services
        LLM externes sécurisés par un proxy de pseudonymisation
        (NER + Regex), avec journalisation et contrôle d'accès ;
  \item \textbf{Approche On-Premise (Souveraine)} : déploiement local
        de modèles open-source (Llama~3, Mistral) sur l'infrastructure
        de la CMR, garantissant la maîtrise totale des données sensibles.
\end{itemize}



\section{Recommandations opérationnelles}

Enfin, ce stage avait pour finalité de produire :

\begin{itemize}
  \item Une \textbf{feuille de route opérationnelle} priorisée pour
        l'adoption sécurisée des LLM à la CMR ;
  \item Des \textbf{politiques de sécurité} adaptées (pseudonymisation,
        contrôle d'accès, gouvernance de l'IA) ;
  \item Un \textbf{tableau comparatif} des solutions cloud et
        on-premise selon les critères de sécurité, coût et conformité ;
  \item Des \textbf{recommandations d'architecture} argumentées, fondées
        sur l'analyse des avantages et inconvénients de chaque approche.
\end{itemize}

Ce stage a couvert une dimension à la fois \textbf{normative},
\textbf{technique} et \textbf{stratégique}, offrant une vision globale
et cohérente de la sécurisation des LLM au sein d'un établissement
public marocain soumis à de fortes exigences de confidentialité et de
conformité.

% ════════════════════════════════════════════════════════════════════
%  ANNEXE – RAPPORT TECHNIQUE COMPLET
% ════════════════════════════════════════════════════════════════════
\appendix

\titleformat{\chapter}[display]
  {\color{primary}\normalfont\huge\bfseries}
  {\color{accent}Annexe\ \thechapter}
  {10pt}
  {\color{primary}\Huge}
  [{\color{accent}\titlerule[0.6pt]}]

\chapter{Étude Technique : Sécurisation des LLM à la CMR}

\begin{tcolorbox}[cmrbox, title={Note sur cette Annexe}]
Cette annexe constitue l’étude technique complet réaliser dans le cadre
du stage. Elle détaille l'analyse de la problématique, les solutions
technologiques étudiées, les comparaisons effectuées et les recommandations
finales formulées à l'attention de la Direction de la Transformation
Digitale de la CMR.
\end{tcolorbox}

\newpage

% ────────────────────────────────────────────────────────────────────
%  A.1 – INTRODUCTION ET PROBLÉMATIQUE
% ────────────────────────────────────────────────────────────────────
\section{Introduction et Problématique}

\subsection{Contexte}

L'explosion de l'intelligence artificielle générative, notamment avec des
modèles de langage large (LLM) comme ChatGPT, Gemini et Claude, offre un
potentiel de productivité sans précédent pour les administrations et
entreprises. Ces outils permettent l'analyse de documents, l'aide à la
décision, la rédaction automatique et la traduction en quelques secondes.
Au Maroc, comme partout ailleurs, les institutions gouvernementales et
entreprises souhaitent exploiter ces capacités.

Cependant, cette révolution technologique s'accompagne d'un défi majeur :
\textbf{comment utiliser la puissance de ces modèles tout en préservant la
confidentialité des données sensibles ?}

\subsection{La Problématique}

La question centrale de ce projet est la suivante :
\textit{Comment concilier l'usage de modèles ultra-performants (souvent
situés sur des serveurs étrangers ou dans le Cloud public international)
avec l'obligation de confidentialité et la souveraineté des données ?}

Cette tension crée un dilemme pour les organisations stratégiques
(institutions gouvernementales, grandes entreprises nationales,
institutions financières comme la CMR) qui ont besoin de l'IA mais ne
peuvent pas risquer de compromettre leurs secrets et données sensibles.

\subsection{Les Risques Identifiés}

\subsubsection{Le Shadow AI}

Les employés utilisent des outils d'IA gratuits (ChatGPT, Gemini) sans
contrôle de leur employeur, contournant ainsi les politiques de sécurité
informatique. Cela crée une « ombre » impossible à tracer et à auditer,
particulièrement critique pour les institutions gérant des données
sensibles.

\subsubsection{L'utilisation des données pour le ré-entraînement}

Certains fournisseurs d'IA peuvent réutiliser les données envoyées par
les utilisateurs pour améliorer leurs modèles, selon les paramètres de
confidentialité et les contrats en place. Une donnée confidentielle
envoyée aujourd'hui pourrait, dans le pire des cas, être exposée
indirectement à travers des comportements non souhaités du modèle ou
via une mauvaise configuration des options de conservation.

\subsubsection{La fuite de Propriété Intellectuelle ou de données d'État}

Pour les administrations gouvernementales ou les entreprises, la fuite
involontaire de données sensibles, de stratégies nationales ou
d'informations personnelles sur les citoyens peut avoir des conséquences
graves, notamment pour les institutions gérant des données personnelles
massives.

\newpage

% ────────────────────────────────────────────────────────────────────
%  A.2 – ANALYSE DU BESOIN
% ────────────────────────────────────────────────────────────────────
\section{Analyse du Besoin}

\subsection{Besoins Métier}

Les organisations souhaitent exploiter les LLM pour :

\begin{itemize}
  \item \textbf{Analyse de documents} : résumer rapidement des contrats,
        décisions administratives ou rapports volumineux ;
  \item \textbf{Aide à la décision} : traiter des données pour fournir
        des recommandations stratégiques ;
  \item \textbf{Rédaction automatique} : générer des rapports
        administratifs, correspondances ou documents officiels ;
  \item \textbf{Traduction} : traduire des documents entre plusieurs
        langues.
\end{itemize}

\subsection{Besoins Sécuritaires}

Pour que l'utilisation des LLM soit acceptable, quatre garanties
fondamentales doivent être assurées.

\subsubsection{Pseudonymisation / Tokenisation}

Remplacement des PII (\textit{Personally Identifiable Information}) tels
que les noms des citoyens, adresses, numéros d'identification, numéros de
carte bancaire, avant toute transmission à un service externe, via des
\textit{tokens} réversibles (ex : \texttt{[PERSONNE\_1]}). Pour une
institution comme la CMR, gérant les dossiers de millions de retraités,
cette mesure est critique.

\begin{tcolorbox}[cmrbox, title={Remarque importante}]
La pseudonymisation est \textbf{réversible} (via la table de
correspondance) et ne doit pas être confondue avec l'anonymisation.
Elle réduit le risque d'exposition, mais les données restent considérées
comme \textbf{données personnelles} du point de vue conformité.
\end{tcolorbox}

\subsubsection{Traçabilité}

Maintenir un journal d'audit complet : qui a envoyé quoi, quand, et
quelle réponse a été reçue. Cela est essentiel pour les investigations
légales ou de conformité, particulièrement pour les institutions gérant
des données massives.

\subsubsection{Non-mémorisation}

Garantir contractuellement que le fournisseur d'IA n'enregistre pas les
données sensibles dans ses serveurs et ne les utilise pas pour
ré-entraîner ses modèles. Cela passe par des clauses
\textit{no retention/no training} et une configuration explicite de
non-conservation côté fournisseur.

\subsubsection{Chiffrement en transit et au repos}

Les données sensibles doivent être chiffrées lors de leur transmission
et, si stockage local il y a, au repos, en conformité avec les standards
de sécurité nationaux.

\subsection{Besoins Juridiques : Cadre Légal Marocain}

Le cadre juridique marocain impose des obligations strictes en matière de
protection des données et de cybersécurité, à travers deux textes
fondamentaux qui encadrent directement l'usage des LLM dans les
institutions publiques.

\subsubsection{Loi n°~09-08 relative à la Protection des Personnes
Physiques à l'égard du Traitement des Données à Caractère Personnel}

Promulguée en 2009, la loi 09-08 constitue le socle juridique de la
protection des données personnelles au Maroc. Elle s'applique à tout
projet impliquant des LLM manipulant des données de citoyens.

\textbf{Obligations clés pour le projet :}
\begin{itemize}
  \item \textbf{Consentement et finalité :} tout traitement de données
        personnelles doit être fondé sur une base légitime et une
        finalité déterminée. L'envoi de données à un LLM externe
        constitue un traitement qui doit être encadré et justifié ;
  \item \textbf{Proportionnalité et minimisation :} seules les données
        strictement nécessaires au traitement doivent être transmises.
        La pseudonymisation avant envoi à un service externe répond
        directement à cette exigence ;
  \item \textbf{Transfert hors du territoire national :} le transfert de
        données personnelles vers un pays étranger est encadré. L'envoi
        de données brutes à des serveurs Cloud situés à l'étranger sans
        mesures adéquates constitue un risque majeur de non-conformité ;
  \item \textbf{Sécurité des traitements :} le responsable du traitement
        doit mettre en œuvre des mesures techniques et organisationnelles
        appropriées contre la perte, l'altération ou la diffusion non
        autorisée ;
  \item \textbf{Droits des personnes :} les citoyens conservent un droit
        d'accès, de rectification et d'opposition. L'institution doit
        pouvoir retracer les traitements et justifier les finalités.
\end{itemize}

\textbf{Organe de contrôle :} la \textbf{CNDP} (Commission Nationale de
contrôle de la protection des Données à caractère Personnel) est
l'autorité chargée de veiller au respect de cette loi. Tout traitement
de données personnelles impliquant des LLM doit être encadré et, selon
le cas, déclaré ou autorisé.

\textbf{Sanctions :} la loi 09-08 prévoit des sanctions pénales et
financières en cas d'infractions (cf. texte officiel).

\subsubsection{Loi n°~05-20 relative à la Cybersécurité}

Promulguée en 2020, la loi 05-20 établit le cadre national de
cybersécurité et s'applique aux infrastructures d'information sensibles,
dont font partie les systèmes d'information de la CMR en tant
qu'institution publique gérant des données critiques.

\textbf{Obligations clés pour le projet :}
\begin{itemize}
  \item \textbf{Classification des systèmes d'information :} les
        systèmes des entités publiques doivent être classifiés selon
        leur niveau de sensibilité. Un système utilisant des LLM pour
        traiter des données de retraités relève d'un niveau de
        sensibilité élevé ;
  \item \textbf{Obligation de sécurisation :} les entités concernées
        doivent appliquer les directives de la \textbf{DGSSI} (Direction
        Générale de la Sécurité des Systèmes d'Information), incluant
        le chiffrement, le contrôle d'accès et l'audit ;
  \item \textbf{Notification des incidents :} en cas d'incident de
        sécurité (y compris exposition potentielle de données),
        l'institution peut être tenue de notifier l'autorité compétente ;
  \item \textbf{Audit et conformité :} les systèmes d'information
        sensibles doivent faire l'objet d'audits réguliers. Tout
        système intégrant des LLM devra être inclus dans le périmètre
        d'audit ;
  \item \textbf{Hébergement souverain :} les données sensibles des
        entités publiques doivent être hébergées sur des infrastructures
        conformes aux exigences nationales, ce qui renforce l'intérêt
        d'un déploiement local ou sur Cloud souverain conforme.
\end{itemize}

\textbf{Organe de contrôle :} la \textbf{DGSSI}, rattachée à
l'Administration de la Défense Nationale, est l'autorité nationale en
matière de cybersécurité. Elle émet les référentiels de sécurité que les
institutions publiques doivent respecter.

\subsubsection{Implications Combinées pour le Projet}

La combinaison des lois 09-08 et 05-20 crée un cadre contraignant mais
clair :

\begin{itemize}
  \item \textbf{Toute donnée personnelle} envoyée à un LLM externe sans
        mesures adaptées (minimisation, pseudonymisation, base légale,
        encadrement contractuel) constitue un risque de non-conformité ;
  \item \textbf{Tout système d'IA} traitant des données sensibles d'une
        entité publique doit respecter les directives de la DGSSI
        conformément à la loi 05-20 ;
  \item \textbf{L'hébergement local} ou sur Cloud souverain conforme est
        fortement recommandé pour les données critiques des institutions
        publiques ;
  \item \textbf{L'auditabilité} de toutes les interactions avec les LLM
        est une exigence de conformité, pas une option technique.
\end{itemize}

Ces contraintes imposent une architecture qui combine confidentialité,
contrôle d'accès et auditabilité.

\newpage

% ────────────────────────────────────────────────────────────────────
%  A.3 – L'APPROCHE PASSERELLE
% ────────────────────────────────────────────────────────────────────
\section{Solutions Technologiques -- L'Approche Passerelle}

\subsection{Le Proxy de Pseudonymisation (Middleware)}

L'idée centrale est d'insérer une couche logicielle intermédiaire entre
l'utilisateur et le service d'IA en ligne.

\subsubsection{Architecture générale}

\begin{enumerate}
  \item L'employé tape sa requête contenant des données sensibles.
  \item Le \textbf{middleware de pseudonymisation} intercepte la requête.
  \item Les PII sont détectées et remplacées par des \textit{tokens}
        (ex : \texttt{[PERSONNE\_1]}, \texttt{[ADRESSE\_A]}).
  \item La requête « nettoyée » est envoyée au service d'IA en ligne
        (ChatGPT, Gemini, etc.).
  \item La réponse est reçue et les tokens sont \textbf{ré-identifiés
        de manière contrôlée} via une table de correspondance interne,
        sans jamais que le fournisseur n'ait vu les vraies données.
  \item Le document final, pseudonymisé du point de vue du fournisseur,
        est présenté à l'utilisateur.
\end{enumerate}

\subsubsection{Détection des Entités}

La détection repose sur deux techniques complémentaires :

\begin{itemize}
  \item \textbf{Expressions régulières (Regex) :} détection simple des
        e-mails, numéros de téléphone, numéros d'identification, etc.
  \item \textbf{Modèles NER :} utilisation de bibliothèques comme
        \textit{Microsoft Presidio} ou \textit{SpaCy} pour identifier
        intelligemment les noms de personnes, les lieux et les
        organisations, même s'ils ne suivent pas un motif prédéfini.
\end{itemize}

\paragraph{Risques de Sécurité du Pipeline NER}

Le moteur NER constitue lui-même un composant sensible de l'architecture
de sécurité. S'il est compromis ou défaillant, c'est l'ensemble de la
chaîne de protection qui s'effondre. Les risques principaux sont :

\begin{itemize}
  \item \textbf{Faux négatifs (PII non détectées) :} si le modèle NER
        ne reconnaît pas une entité sensible (nom rare, format d'adresse
        inhabituel, numéro d'identification non standard), cette donnée
        sera envoyée en clair au fournisseur Cloud --- risque critique ;
  \item \textbf{Attaques adversariales :} un utilisateur malveillant
        pourrait volontairement formater ses requêtes pour contourner la
        détection NER (insertion de caractères invisibles,
        translittération, encodage alternatif) ;
  \item \textbf{Empoisonnement du modèle NER :} si le modèle NER est
        ré-entraîné sur des données internes, un acteur malveillant
        pourrait injecter des données d'entraînement corrompues pour
        affaiblir la détection de certaines catégories de PII ;
  \item \textbf{Fuite via la table de correspondance :} la table qui
        associe les tokens aux données réelles constitue un actif
        hautement sensible. Si cette table est compromise, toutes les
        données originales sont exposées.
\end{itemize}

\paragraph{Mesures de Sécurisation du NER}

Pour garantir la fiabilité du pipeline NER, les mesures suivantes doivent
être implémentées :

\begin{itemize}
  \item \textbf{Double couche de détection :} combiner Regex (détection
        par motifs) et NER sémantique (détection par contexte). Une PII
        manquée par l'un sera attrapée par l'autre ;
  \item \textbf{Liste blanche de sortie :} appliquer un filtre vérifiant
        que la requête pseudonymisée ne contient aucun motif ressemblant
        à un PII connu avant envoi au Cloud ;
  \item \textbf{Chiffrement de la table de correspondance :} chiffrement
        AES-256 au repos, clés stockées dans un HSM ou coffre-fort de
        secrets (ex : HashiCorp Vault) ;
  \item \textbf{Durée de vie limitée des mappings :} tables de
        correspondance éphémères, détruites automatiquement après
        ré-identification. Aucune persistance à long terme ;
  \item \textbf{Tests adversariaux réguliers :} tests de pénétration
        spécifiques au NER --- soumettre des requêtes volontairement
        obfusquées pour valider la résistance aux contournements ;
  \item \textbf{Exécution locale isolée :} le moteur NER doit
        impérativement s'exécuter sur les serveurs internes de
        l'institution, jamais sur un service Cloud ;
  \item \textbf{Audit et versionnage :} chaque version du modèle NER
        déployée doit être auditée, et les résultats enregistrés pour
        analyse post-incident.
\end{itemize}

\begin{tcolorbox}[cmrbox, title={Limite structurelle de l'approche Passerelle}]
Malgré ces mesures, la détection automatique ne peut pas garantir un
taux de capture de 100\% dans tous les cas (formats inattendus,
obfuscation). C'est pourquoi l'approche Passerelle doit rester une
\textbf{solution de compromis}, et non un socle unique pour une
institution critique.
\end{tcolorbox}

\subsection{Avantages de l'Approche Passerelle}

\begin{itemize}
  \item Permet d'utiliser les \textbf{meilleurs modèles cloud du marché}
        sans envoyer de données personnelles en clair ;
  \item Coût plus faible que l'on-prem (pas de GPU), mais nécessite une
        sécurité applicative sérieuse ;
  \item Rapidité de déploiement ;
  \item Les modèles continuent de bénéficier des mises à jour du
        fournisseur.
\end{itemize}

\newpage

% ────────────────────────────────────────────────────────────────────
%  A.4 – L'APPROCHE SOUVERAINE
% ────────────────────────────────────────────────────────────────────
\section{Solutions Technologiques -- L'Approche Souveraine et Sécurité
Avancée}

\subsection{Le Déploiement Local (On-Premise)}

Au lieu de faire confiance à un fournisseur externe, l'organisation
héberge son propre modèle d'IA sur ses serveurs internes.

\subsubsection{Modèles Open Source Recommandés}

\begin{itemize}
  \item \textbf{Mistral 7B/8x7B :} très performant et léger ;
  \item \textbf{Llama~3 (70B) :} capacités comparables à des modèles
        généralistes pour de nombreux cas d'usage métier ;
  \item \textbf{Phi~3 :} optimisé pour les appareils légers.
\end{itemize}

\subsubsection{Hébergement Sécurisé}

Dans ce rapport, « souverain » signifie : hébergement interne ou cloud
localisé sur le territoire national avec garanties légales et
contractuelles.

\begin{itemize}
  \item \textbf{Serveurs internes :} les institutions hébergent les
        modèles sur leurs propres serveurs, sous contrôle total de
        l'institution ;
  \item \textbf{Cloud Souverain :} des prestataires certifiés offrent un
        Cloud situé sur le territoire national avec des garanties légales
        conformes aux lois de protection des données.
\end{itemize}

\subsection{Le RAG Local : Un Dispositif de Sécurité}

Le RAG (\textit{Retrieval Augmented Generation}) n'est pas seulement une
technique d'amélioration des performances. Dans ce contexte, il constitue
un \textbf{mécanisme de sécurité fondamental} qui applique le principe du
\textbf{moindre privilège informationnel} : l'IA ne reçoit jamais
l'ensemble des données, mais uniquement les extraits strictement
nécessaires pour répondre à une question précise.

\begin{tcolorbox}[cmrbox, title={Pourquoi le RAG est un Outil de Sécurité}]
Sans RAG, pour qu'un LLM connaisse les données de l'organisation, il
faudrait soit les inclure intégralement dans chaque requête (risque
d'exposition massive), soit ré-entraîner le modèle sur ces données
(le modèle « mémorise » alors les informations sensibles de manière
durable). Ces deux approches sont des \textbf{failles de sécurité
structurelles}.
\end{tcolorbox}

\subsubsection{Gains de Sécurité du RAG}

\begin{enumerate}
  \item \textbf{Isolation des données sensibles :} documents dans une
        base vectorielle locale, sur les serveurs de l'institution. Le
        LLM ne reçoit que des fragments contextuels pertinents, jamais
        l'intégralité d'un dossier ;
  \item \textbf{Élimination du risque de mémorisation :} contrairement
        au fine-tuning, le RAG ne modifie pas les poids du modèle. Les
        données ne sont jamais « apprises » par l'IA : elles sont
        consultées ponctuellement puis oubliées à la fin de la session ;
  \item \textbf{Contrôle d'accès documentaire :} la base vectorielle
        peut être segmentée par domaine métier (Retraites, RH, Finance)
        avec RBAC ;
  \item \textbf{Traçabilité des sources :} chaque réponse peut citer ses
        sources exactes (document, page, paragraphe), permettant un audit
        précis des contenus utilisés ;
  \item \textbf{Réduction des hallucinations :} en ancrant les réponses
        dans des documents fournis, le RAG réduit fortement ce risque
        opérationnel ;
  \item \textbf{Surface d'attaque minimisée :} sans RAG, un attaquant
        pourrait tenter d'extraire des informations mémorisées. Avec le
        RAG, les données ne sont pas stockées dans les paramètres du
        modèle.
\end{enumerate}

\subsubsection{Fonctionnement Technique}

\begin{enumerate}
  \item \textbf{Vectorisation (Embeddings) :} les documents confidentiels
        (PDFs, rapports, emails) sont convertis en vecteurs dans une base
        locale sécurisée, segmentée par domaine métier et soumise au
        contrôle d'accès RBAC ;
  \item \textbf{Recherche Sémantique :} lorsqu'une requête est lancée,
        le moteur extrait uniquement les paragraphes pertinents de la
        base documentaire, limités au périmètre autorisé de l'utilisateur ;
  \item \textbf{Génération Contrainte :} le modèle reçoit ces extraits
        avec une instruction stricte : \textit{``Réponds uniquement en
        te basant sur ces documents.''}.
\end{enumerate}

\subsubsection{Validation du RAG Sécurisé}

Pour garantir que le RAG respecte les contrôles d'accès, une validation
empirique est recommandée :

\begin{itemize}
  \item \textbf{Test d'isolation documentaire :} un utilisateur RH tente
        d'accéder à des documents Finance via requête. Le RAG doit
        refuser (\texttt{UNAUTHORIZED\_DOMAIN\_ACCESS}) ;
  \item \textbf{Test du NER :} un pipeline de 50 requêtes réelles
        (anonymisées) avec PII doit être exécuté. Taux de capture
        attendu : \(\geq 98\%\) pour les PII standards.
\end{itemize}

Ces tests doivent être répétés trimestriellement post-déploiement.

\subsection{Besoins de Sécurité Interne et Contrôle d'Accès}

\subsubsection{Problématique d'Accès Hiérarchique}

Un manager ou une personne hiérarchiquement supérieure ne doit pas
nécessairement avoir accès à toutes les données de l'organisation.

\textit{Exemple :}
\begin{itemize}
  \item Un agent Retraites partage des données sensibles de retraités
        avec l'IA pour une analyse ;
  \item Un manager RH ne doit pas voir ces données des retraités, même
        s'il est hiérarchiquement supérieur à l'agent ;
  \item Un administrateur système ne doit pas pouvoir déchiffrer les
        données financières sensibles de la CMR.
\end{itemize}

Cette séparation par domaine métier, indépendante de la hiérarchie, est
essentielle pour respecter le principe du \og besoin de connaître \fg{}
(\textit{Need-to-Know Principle}).

\subsubsection{RBAC -- Fondation}

Principe : chaque employé se voit attribuer des rôles spécifiques basés
sur son poste.

\textit{Rôles typiques à la CMR :}
\begin{itemize}
  \item \texttt{Agent\_Retraites} : accès uniquement à l'IA Retraites ;
  \item \texttt{Agent\_Paie} : accès uniquement à l'IA Paie ;
  \item \texttt{Manager\_RH} : accès à l'IA RH et données RH seulement ;
  \item \texttt{Admin\_Systeme} : maintenance uniquement, pas d'accès
        données métier ;
  \item \texttt{Responsable\_Conformite} : consultation des audit logs,
        pas des données opérationnelles.
\end{itemize}

\subsubsection{Audit et Logging Granulaire en Temps Réel}

Principe : chaque accès à l'IA et aux données est enregistré en détail.

\textit{Données enregistrées :}
\begin{itemize}
  \item User ID (identité de la personne) ;
  \item Timestamp (date et heure exacte) ;
  \item User Role (son rôle dans le système) ;
  \item IA Domain Accessed (IA Retraites, Finances, RH, etc.) ;
  \item Action (send request, retrieve result, etc.) ;
  \item Number of PII Detected (nombre de données sensibles traitées) ;
  \item Status (\texttt{SUCCESS} si autorisé, \texttt{BLOCKED} si non
        autorisé).
\end{itemize}

\textbf{Alertes automatiques :} déclenchées si un utilisateur tente
d'accéder à une IA non autorisée.\\
\textbf{Dashboard :} un tableau de bord permet de visualiser les accès,
les violations et les tendances.

\subsubsection{Charte d'Utilisation et Formations Régulières}

\begin{itemize}
  \item Session de 30 minutes chaque trimestre sur les risques et la
        conformité ;
  \item Quiz de certification obligatoire avant premier accès à l'IA ;
  \item Refresher annuel pour tous les utilisateurs.
\end{itemize}

\subsection{Architecture Multi-Agents à Politique Fixe}

Afin d'éviter les limites et risques liés à l'injection dynamique de
\textit{system prompts}, nous proposons une \textbf{architecture
multi-agents} composée d'agents IA distincts, chacun avec :

\begin{itemize}
  \item un \textbf{rôle métier} clair ;
  \item un \textbf{périmètre fonctionnel} strict (outils, bases, RAG) ;
  \item une \textbf{politique d'accès fixe} (avec ou sans
        authentification) ;
  \item un \textbf{prompt système statique}.
\end{itemize}

Cette approche améliore la conformité et la sécurité, car chaque agent
devient un composant stable, contrôlable et facilement auditable :
l'organisation sait précisément \textit{qui peut accéder à quoi}, via
\textit{quel agent}, et sous \textit{quelles conditions}.

\subsubsection{Les 3 Types d'Agents}

\paragraph{Type 1 -- Agent Sensible (Authentification Obligatoire)}
Cet agent traite des requêtes qui \textbf{nécessitent un accès à des
données internes confidentielles} (paies, pensions, budgets, dossiers RH,
états financiers). Il \textbf{n'existe pas} de mode sans identification
pour ce type d'agent.

\textit{Exemples :}
\begin{itemize}
  \item Agent Finance : consultation des paiements, budgets,
        financements, rapports comptables ;
  \item Agent RH : dossiers du personnel, documents internes, décisions
        non publiques ;
  \item Agent Retraites : dossiers de retraités, montants, historiques,
        matricules.
\end{itemize}

\textbf{Politique de sécurité :}
\begin{itemize}
  \item Authentification obligatoire dès l'accès à l'agent ;
  \item Vérification \textbf{RBAC} systématique ;
  \item Accès documentaire interne via \textbf{RAG} (segments autorisés
        uniquement) ;
  \item Logs complets et alertes en cas d'accès non autorisé.
\end{itemize}

\paragraph{Type 2 -- Agent Public (Sans Authentification)}
Cet agent répond à des besoins basés sur \textbf{des informations
publiques} et ne doit \textbf{jamais} accéder aux données internes.

\textit{Exemple :}
\begin{itemize}
  \item Agent Juridique Public : lois, textes officiels publiés,
        réglementations publiques, jurisprudence accessible.
\end{itemize}

\textbf{Politique de sécurité :}
\begin{itemize}
  \item Aucun login requis ;
  \item Aucun accès aux bases internes (pas de RAG interne) ;
  \item Si un RAG existe, il est limité à un \textbf{corpus public
        validé} ;
  \item Réponses bornées avec citation des sources publiques (BO, textes
        officiels, référentiels) ;
  \item Journalisation légère (statistiques d'usage), sans traitement de
        données personnelles.
\end{itemize}

\paragraph{Type 3 -- Agent Mixte (Libre si Non Sensible, Bascule si
Sensible)}
Cet agent est accessible sans identification pour les requêtes non
sensibles (rédaction neutre, méthodologie, traduction de texte public,
assistance bureautique). Si la requête est détectée comme sensible, le
système \textbf{bascule automatiquement} vers une authentification.

\textbf{Politique de sécurité :}
\begin{itemize}
  \item Accès libre par défaut, \textbf{sans accès interne} ;
  \item Analyse de sensibilité des requêtes (NER + mots-clés + règles) ;
  \item En cas de doute, politique conservative : la requête est traitée
        comme sensible ;
  \item Si sensible : interruption $\rightarrow$ authentification
        $\rightarrow$ vérification RBAC $\rightarrow$ RAG interne
        autorisé $\rightarrow$ logs complets.
\end{itemize}

\subsubsection{Mécanisme de Détection et Bascule (Agent Mixte)}

Le basculement vers le mode sécurisé est déclenché par :

\begin{itemize}
  \item \textbf{Détection de PII} (noms, adresses, identifiants,
        matricules, numéros) ;
  \item \textbf{Mots-clés sensibles} (pension, paie, budget, dossier,
        matricule, financement, etc.) ;
  \item \textbf{Demande explicite d'accès interne} (analyse d'un
        dossier ou d'un document interne via RAG) ;
  \item \textbf{Principe de prudence :} en cas d'incertitude sur la
        sensibilité, le système applique une politique conservative et
        déclenche l'authentification.
\end{itemize}

Lorsque la requête est jugée sensible :

\begin{enumerate}
  \item le système \textbf{bloque} l'exécution avant tout traitement ;
  \item demande une \textbf{authentification} ;
  \item vérifie les \textbf{habilitations RBAC} ;
  \item route la requête vers l'agent sensible correspondant (Type~1)
        ou active les contrôles sécurisés ;
  \item enregistre l'ensemble de la transaction dans les
        \textbf{audit logs}.
\end{enumerate}

\subsection{Serveurs MCP et Contrôle d'Accès Inter-Domaines}

Le MCP (\textit{Model Context Protocol}) désigne ici la couche
\textbf{d'orchestration des outils} (\textit{tools}) et le point
d'intégration entre le LLM et les ressources (authentification, RAG,
bases internes, journaux, etc.). Dans cette architecture, le LLM analyse
la requête utilisateur, propose un \textbf{plan d'action} et sélectionne
les \textbf{tools} exposés via le MCP
(ex : \texttt{auth\_challenge()}, \texttt{rag\_retrieve()},
\texttt{audit\_log()}).

Le MCP joue le rôle de \textbf{Policy Enforcement Point (PEP)} :

\begin{itemize}
  \item il \textbf{valide} chaque appel d'outil avant exécution ;
  \item il \textbf{impose} l'authentification lorsque requis ;
  \item il \textbf{vérifie} les habilitations (RBAC/ABAC) pour tout
        accès à un domaine ou une ressource ;
  \item il \textbf{isole} strictement les domaines (un agent public ne
        peut jamais appeler un tool RAG interne) ;
  \item il \textbf{journalise} les décisions pour audit et conformité.
\end{itemize}

\subsubsection{Workflow de Contrôle (Tool-Gated)}

\begin{enumerate}
  \item \textbf{Analyse LLM :} le modèle interprète la requête et tente
        d'utiliser les tools nécessaires ;
  \item \textbf{Garde-fou MCP (PEP) :} à chaque tool call, le MCP
        vérifie le contexte (session, identité, rôle, domaine demandé) ;
  \item \textbf{Authentification si nécessaire :} si la requête implique
        un accès sensible, le MCP impose un défi via
        \texttt{auth\_challenge()} ;
  \item \textbf{Autorisation :} après authentification, le MCP autorise
        ou refuse les tools sensibles
        (ex : \texttt{rag\_retrieve(domain=RH)}) ;
  \item \textbf{Cloisonnement inter-domaines :} toute tentative de
        changement de domaine déclenche une nouvelle vérification
        d'habilitation ;
  \item \textbf{Audit :} chaque décision est loggée (tool demandé,
        domaine, rôle, statut, horodatage).
\end{enumerate}

\subsubsection{Validation (Tests de Politique Côté Tools)}

Avant déploiement, valider que la politique est \textbf{enforcée
indépendamment du prompt} :

\begin{itemize}
  \item 100 requêtes non sensibles : accès autorisé sans authentification
        (pas de tools internes) ;
  \item 100 requêtes sensibles : accès \textbf{bloqué} tant que
        \texttt{auth\_challenge()} n'est pas validé ;
  \item Tests inter-domaines : un utilisateur RH ne doit jamais accéder
        au domaine Finance (tool \texttt{rag\_retrieve} refusé) ;
  \item Tests prompt-injection : tentatives de forcer l'accès interne
        via instructions malveillantes $\rightarrow$ refus côté MCP.
\end{itemize}

\newpage

% ────────────────────────────────────────────────────────────────────
%  A.5 – COMPARAISONS ET INCONVÉNIENTS
% ────────────────────────────────────────────────────────────────────
\section{Comparaisons et Inconvénients}

\subsection{Tableau Comparatif}

\begin{table}[H]
\centering
\renewcommand{\arraystretch}{1.3}
\setlength{\tabcolsep}{6pt}
\begin{tabular}{|p{4.0cm}|p{5.5cm}|p{5.5cm}|}
\hline
\rowcolor{primary}
\textbf{\color{white}Critère}
  & \textbf{\color{white}Cloud Public (+ pseudonymisation)}
  & \textbf{\color{white}On-Premise (déploiement local)} \\
\hline
\textbf{Confidentialité}
  & Moyenne : dépend du pipeline (Regex/NER), du stockage des logs et
    des garanties du fournisseur
  & Très élevée : les données restent dans le réseau interne \\
\hline
\textbf{Souveraineté}
  & Faible : dépendance à un prestataire externe
  & Complète : contrôle total par l'institution \\
\hline
\textbf{Conformité 09-08 / 05-20}
  & Partielle : transfert + risque résiduel même pseudonymisé
  & Forte : réduction structurelle du transfert hors territoire \\
\hline
\textbf{Risque ré-entraînement}
  & Modéré : dépend des options et clauses \textit{no training /
    no retention}
  & Nul : aucune donnée ne sort du SI \\
\hline
\textbf{Coût infrastructure}
  & Bas : paiement à l'usage, pas de GPU interne
  & Élevé : GPU, énergie, refroidissement, maintenance \\
\hline
\textbf{Coût d'exploitation}
  & Moyen : sécurité applicative (proxy, NER, audit) + coûts API
  & Élevé : MLOps, DevOps, supervision, scalabilité \\
\hline
\textbf{Performance modèle}
  & Excellente : modèles premium cloud
  & Bonne : open source + RAG (selon dimensionnement) \\
\hline
\textbf{Latence}
  & 1--5 secondes
  & \(<\) 1 seconde (réseau interne) \\
\hline
\textbf{Mises à jour}
  & Automatiques (côté fournisseur)
  & Manuelles (cycle interne) \\
\hline
\textbf{Surface d'attaque}
  & Pipeline pseudo + API externe + erreurs de configuration
  & Accès internes + gestion des secrets + RBAC \\
\hline
\end{tabular}
\caption{Comparaison entre Cloud public avec pseudonymisation et
déploiement on-premise}
\label{tab:comparaison_annexe}
\end{table}

\subsection{Inconvénients Majeurs}

\subsubsection{Perte de qualité (Approche Passerelle)}

La pseudonymisation peut introduire des artefacts textuels :

\begin{quote}
\textit{Texte original :} « Mohamed Bennani, résident à Casablanca,
retraité depuis 2020 avec une pension de 8~500~DH. »

\textit{Texte pseudonymisé :} « [PER\_1] résident à [VILLE\_1],
retraité depuis [ANNEE\_1] avec une pension de [MONTANT\_1]. »
\end{quote}

Si l'IA doit générer une phrase grammaticale, les tokens peuvent parfois
casser la cohérence. Une IA moins robuste produira une réponse maladroite.

\subsubsection{Coût d'infrastructure (Approche Locale)}

Maintenir un cluster de serveurs GPU pour exécuter des LLM coûte cher :

\begin{itemize}
  \item Une GPU NVIDIA A100 : \(\sim\) 110~000~DH ;
  \item Électricité et refroidissement : plusieurs milliers de DH/mois ;
  \item Peu adapté aux petites et moyennes structures.
\end{itemize}

\subsubsection{Complexité de maintenance}

L'approche locale nécessite une expertise pointue :

\begin{itemize}
  \item Experts en Machine Learning ;
  \item Experts en cybersécurité et conformité ;
  \item Ingénieurs DevOps/infrastructure ;
  \item Suivi des mises à jour des modèles et des dépendances.
\end{itemize}

\subsubsection{Limitation des modèles open source}

Les meilleurs modèles open source restent parfois en retrait face aux
meilleurs modèles cloud sur certaines tâches très complexes (raisonnement
avancé, code très élaboré), selon les cas d'usage et le dimensionnement.

\newpage

% ────────────────────────────────────────────────────────────────────
%  A.6 – RECOMMANDATIONS ET CONCLUSION
% ────────────────────────────────────────────────────────────────────
\section{Recommandations et Conclusion}

\subsection{Stratégie Recommandée : L'Approche Souveraine}

Après analyse des solutions technologiques, du cadre juridique marocain
et des exigences spécifiques de la CMR, nous recommandons l'adoption de
\textbf{l'Approche Souveraine} comme stratégie principale de sécurisation
de l'usage des LLM.

Cette approche combine le déploiement local de modèles open source, une
architecture RAG sécurisée, et un \textbf{système d'agents IA à politiques
d'accès fixes} (Public / Mixte / Sensible) décrit dans ce rapport.

\subsection{Pourquoi l'Approche Souveraine et pas la Passerelle}

\subsubsection{L'Exigence Légale est Claire}

\begin{itemize}
  \item \textbf{Loi 09-08 :} le transfert de données personnelles hors
        du territoire est encadré. Même pseudonymisées, les requêtes
        envoyées à des serveurs étrangers peuvent contenir du contexte
        métier sensible qui doit être protégé et justifié ;
  \item \textbf{Loi 05-20 :} les entités publiques doivent appliquer les
        exigences nationales de cybersécurité (DGSSI), ce qui renforce
        la préférence pour des solutions maîtrisées ;
  \item \textbf{Risque résiduel :} même avec des clauses
        \textit{no retention} et des options de confidentialité,
        l'institution reste responsable en cas d'incident. Un déploiement
        interne réduit structurellement la dépendance à un tiers.
\end{itemize}

\subsubsection{Sécurité Structurelle vs. Sécurité Dépendante}

L'approche Passerelle repose sur une chaîne de confiance fragile :

\begin{itemize}
  \item Confiance dans le moteur NER pour détecter \textbf{toutes} les
        PII (difficile à garantir) ;
  \item Confiance dans le fournisseur Cloud pour respecter ses
        engagements contractuels et ses options de conservation ;
  \item Confiance dans la sécurisation du transport et des systèmes
        intermédiaires.
\end{itemize}

L'approche interne \textbf{réduit ces dépendances} : les données ne
quittent pas le périmètre de l'institution.

\subsubsection{Maturité des Modèles Open Source}

L'argument historique en faveur du Cloud se réduit :

\begin{itemize}
  \item Des modèles open source modernes offrent des performances
        adaptées à la majorité des tâches métier (résumé, rédaction,
        analyse documentaire) ;
  \item Pour les tâches courantes de la CMR, l'écart peut être
        \textbf{non bloquant} si l'architecture est renforcée par le RAG
        et des contrôles de sécurité.
\end{itemize}

\subsection{Architecture Technique Recommandée}

\begin{enumerate}
  \item \textbf{Serveurs GPU internes} ou \textbf{Cloud Souverain
        conforme} hébergeant des LLM open source ;
  \item \textbf{Base vectorielle locale} (ChromaDB, Weaviate) pour le
        RAG, segmentée par domaine métier avec RBAC ;
  \item \textbf{Portail d'accès unifié} basé sur trois types d'agents :
  \begin{itemize}
    \item \textbf{Agent Sensible} (authentification obligatoire) ;
    \item \textbf{Agent Public} (sans authentification, sources publiques
          uniquement) ;
    \item \textbf{Agent Mixte} (usage libre si non sensible, bascule
          vers authentification si nécessaire).
  \end{itemize}
\end{enumerate}

\textbf{Couches de sécurité intégrées :}
\begin{itemize}
  \item Agents à prompts statiques et politiques d'accès fixes ;
  \item Détection automatique de sensibilité (NER + règles + mots-clés) ;
  \item RAG sécurisé avec segmentation documentaire et traçabilité des
        sources ;
  \item RBAC pour le contrôle d'accès basé sur les rôles métier ;
  \item MCP pour le cloisonnement inter-domaines et la vérification des
        habilitations ;
  \item Audit et logging granulaire de toutes les interactions sensibles ;
  \item Chiffrement au repos et en transit pour les données, index et
        journaux.
\end{itemize}

\subsection{Flux de Traitement Complet}

\begin{enumerate}
  \item L'employé accède au portail IA de la CMR. Par défaut, l'Agent
        Public et l'Agent Mixte sont disponibles. L'accès à tout Agent
        Sensible est \textbf{bloqué} tant que l'utilisateur n'est pas
        authentifié.
  \item \textbf{Agent Public :} réponses basées sur un corpus public,
        sans authentification, sans accès interne.
  \item \textbf{Agent Sensible :} authentification obligatoire
        $\rightarrow$ vérification RBAC $\rightarrow$ accès RAG interne
        autorisé $\rightarrow$ réponse avec citations $\rightarrow$
        logging complet.
  \item \textbf{Agent Mixte :} usage libre pour requêtes non sensibles.
        Si une requête est jugée sensible : interruption $\rightarrow$
        authentification $\rightarrow$ vérification RBAC $\rightarrow$
        routage vers l'agent Sensible approprié $\rightarrow$ RAG
        autorisé $\rightarrow$ réponse $\rightarrow$ logging complet.
  \item Si changement de domaine détecté $\rightarrow$ le MCP intervient
        pour vérifier les habilitations avant tout accès à une ressource
        d'un autre domaine.
  \item Toutes les opérations sensibles sont enregistrées pour audit et
        conformité.
\end{enumerate}

\subsection{Gestion des Inconvénients}

\begin{itemize}
  \item \textbf{Coût d'infrastructure :} commencer avec des modèles
        légers, puis monter en puissance. Comparer le coût
        d'infrastructure au coût d'un incident de sécurité (sanctions,
        réputation, impacts opérationnels). Le Cloud Souverain offre une
        alternative au matériel propre, avec coûts mensualisés ;
  \item \textbf{Expertise technique :} partenariat initial + transfert
        de compétences à l'équipe interne. Les outils modernes (Ollama,
        vLLM, LangChain) simplifient la gestion ;
  \item \textbf{Performance des modèles :} le RAG compense l'écart de
        performance brut en ancrant les réponses dans les documents réels
        de l'institution. Pour les tâches métier, un modèle open source
        avec RAG peut surpasser un modèle cloud généraliste sans contexte
        documentaire. Pour certaines tâches très spécifiques, l'approche
        Passerelle peut être utilisée ponctuellement en complément.
\end{itemize}

\subsection{Validation et Assurance Qualité}

Avant déploiement en production :

\begin{enumerate}
  \item \textbf{Benchmark modèles :} Llama~3 70B vs ChatGPT sur 50 cas
        réels --- performance attendue : \(\geq 85\%\) ;
  \item \textbf{Tests de sécurité :} pen testing basique (injection de
        prompt, tentative d'escalade RBAC) ;
  \item \textbf{Conformité 09-08/05-20 :} audit par le service juridique.
\end{enumerate}

Ces validations prennent 4 à 6 semaines et doivent précéder tout
déploiement en production.

\vspace{0.8cm}

% FIX #5 – remplacement du fcolorbox+parbox instable par tcolorbox
\begin{tcolorbox}[
  enhanced,
  colback=white, colframe=primary,
  arc=4pt, boxrule=1.2pt,
  left=10pt, right=10pt, top=10pt, bottom=10pt,
  title={\large\bfseries Note Importante},
  fonttitle=\large\bfseries\color{white},
  coltitle=white,
  attach boxed title to top center={yshift=-2mm},
  boxed title style={colback=primary, rounded corners}
]
Ce rapport recommande l'Approche Souveraine comme stratégie de
sécurisation des LLM pour la CMR. \textbf{Cette recommandation est
conditionnée} à la validation positive des tests décrits aux sections
suivantes :

\vspace{0.3cm}
\begin{itemize}[leftmargin=2cm]
  \item Section A.4 : Tests d'isolation documentaire du RAG
  \item Section A.4 : Tests de détection de sensibilité de l'Agent Mixte
  \item Section A.6 : Benchmark et tests de sécurité pré-déploiement
\end{itemize}

\vspace{0.2cm}
\textbf{Aucun déploiement en production ne doit être effectué} sans que
ces validations aient produit des résultats positifs et documentés.
\end{tcolorbox}

\vspace{0.8cm}

\subsection{Conclusion}

L'Approche Souveraine est la solution la plus adéquate pour la CMR
car elle :

\begin{itemize}
  \item Garantit la \textbf{souveraineté} des données : aucune donnée
        ne quitte le périmètre de l'institution ;
  \item Assure une \textbf{conformité renforcée} avec les lois 09-08
        et 05-20 ;
  \item Réduit structurellement les risques de ré-entraînement externe,
        de fuite et de dépendance à un tiers ;
  \item Offre des \textbf{performances adaptées} grâce au RAG qui ancre
        les réponses dans les documents internes ;
  \item Maximise l'\textbf{adoption} via les agents Public et Mixte pour
        les usages non sensibles, réduisant le Shadow AI ;
  \item Intègre une \textbf{sécurité proportionnelle} : l'authentification
        et les contrôles stricts ne s'activent que lorsque des données
        confidentielles sont en jeu ;
  \item Permet une \textbf{traçabilité} complète des interactions
        sensibles, utile pour l'audit et la conformité.
\end{itemize}

L'approche Passerelle par pseudonymisation, bien que techniquement
intéressante, repose sur une chaîne de confiance externe (efficacité du
NER, conditions du fournisseur Cloud) qui peut être moins compatible avec
le niveau de souveraineté exigé d'une institution publique gérant des
données critiques.

L'Approche Souveraine représente un choix stratégique qui place la CMR
en position de \textbf{contrôle} sur son usage de l'intelligence
artificielle, aujourd'hui et pour les années à venir.

\vspace{1cm}
\noindent\rule{10cm}{0.4pt}

\end{document}